\documentclass{imc-inf} % Use the actual IMC template class

% --- IMC Template Specific Commands ---
% These commands are defined by the imc-inf.cls file
\title{Optimizing ML Experiments: Intelligent Frameworks for Efficiency and Reproducibility}
\subtitle{A Fingerprint-Based LOSO Caching Approach Validated on EEG Sleep Classification}
\thesistype{Bachelor Thesis Proposal} % FIXED: Changed from Bachelor Exposé
\author{Lennart Gorzel}
\supervisor{Prof. Himanshu Buckchash} % FIXED: Correct supervisor
\copyrightyear{2026} % FIXED: Updated year
\submissiondate{15.01.2026} % FIXED: Proposal deadline
\keywords{Machine Learning, Intelligent Caching, Reproducibility, Sleep Stage Classification, EEG, Feature Engineering}

% --- Standard Packages (keep these) ---
\usepackage[utf8]{inputenc}
\usepackage[english]{babel}
\usepackage{svg}
\usepackage{rotating}
\usepackage{afterpage}
\usepackage{amsmath}
\usepackage{amsfonts}
\usepackage{amssymb}
\usepackage{graphicx}
\usepackage{listings}
\usepackage{subcaption}
\usepackage{booktabs} % For professional tables
\usepackage{multirow} % For merged cells in tables
\usepackage{xcolor}


\begin{document}

% --- Front Matter (as per IMC template) ---
\frontmatter
\maketitle

\begin{declarations}
\end{declarations}

% ============================================================================
% UPDATED ABSTRACT - Proposal Style
% ============================================================================

\begin{abstract}
This thesis addresses the computational costs and reproducibility challenges 
that confront modern Machine Learning (ML) experimentation. As ML systems grow 
increasingly complex, researchers face significant overhead from redundant 
computations, inconsistent results across runs, and difficulties resuming 
interrupted experiments. This work proposes a fingerprint-based caching 
framework using SHA-256 configuration hashing to systematically address these 
challenges.

The research develops a deterministic fingerprinting system where a cryptographic 
hash uniquely identifies each experiment configuration.The fingerprint is computed from a canonical JSON representation of the tracked experiment state: random seed, code version, model type and hyperparameters, feature configuration, and held-out subject. The core principle is: 
\textit{if it can change the output, it must change the fingerprint}. This enables automatic cache invalidation for all tracked experiment parameters: when any fingerprinted component changes, the cache key changes and stale results are bypassed. Identical fingerprints therefore indicate the same tracked configuration and a valid cache match.

The framework extends caching to Leave-One-Subject-Out cross-validation by including the held-out subject in the cache key. This prevents reuse of models across different LOSO folds and avoids cross-contamination between train/test splits under the fixed dataset partition used in this thesis.

The framework is validated through a sleep stage classification case study using the Bitbrain Open Access Sleep (BOAS) dataset (128 subjects, approximately 120,000 epochs). 
The feature extraction module computes 149 features per epoch across 6 EEG channels (F3, F4, C3, C4, O1, O2), spanning time-domain statistics, frequency-domain power spectral features, and signal complexity measures. These features are computed once per subject and cached persistently, isolating the caching evaluation to the computationally expensive model training phase.

\begin{center}
\colorbox{yellow!20}{%
\begin{minipage}{0.9\textwidth}
\textbf{Working Draft Note:} 
check speedup claims 
\end{minipage}
}
\end{center}

The contributions of this thesis are: (1) a fingerprint-based LOSO model caching framework that ensures cache correctness through SHA-256 configuration tracking, preventing cross-contamination across folds, (2) quantified efficiency gains of up to 498x speedup for XGBoost and 110x for Random Forest at the highest feature dimensionality (149 features, 128 LOSO folds) and (3) validated reproducibility — all cached models produce identical classification results to cold retraining (accuracy match confirmed for all configurations).
\end{abstract}

\begin{center}
\colorbox{yellow!20}{%
\begin{minipage}{0.9\textwidth}
\textbf{Working Draft Note:} 
rework table of contentns, make sure sraita goal of chapter lenght is achieved
\end{minipage}
}
\end{center}

\addtoToC{Table of Contents}%
\tableofcontents%
\clearpage

\addtoToC{List of Tables}%
\listoftables
\clearpage

\addtoToC{List of Figures}%
\listoffigures
\clearpage

% ============================================================================
% MAIN MATTER
% ============================================================================
\mainmatter



\begin{center}
\colorbox{yellow!20}{%
\begin{minipage}{0.9\textwidth}
\textbf{Working Draft Note:} 
This document currently uses structured lists for clarity during the planning phase. 
These will be converted to prose format in the final thesis submission.
\end{minipage}
}
\end{center}
% ============================================================================
% CHAPTER 1: INTRODUCTION (Keep existing content, just update where needed)
% ============================================================================

\chapter{Introduction}

\section{Introduction and Motivation}
The field of Artificial Intelligence (AI) and Machine Learning (ML) is experiencing unprecedented growth, marked by an explosive increase in both academic publications and industrial applications. This rapid expansion has led to the development of ML systems that present unique challenges compared to traditional software. Their complexity stems primarily from a reliance on dynamic data, evolving model behaviors, and continuous learning paradigms, factors that conventional software management practices are not designed to address. This complexity extends across the entire ML lifecycle, encompassing data pipelines, model deployment, and continuous monitoring.

The sheer volume of research output underscores this challenge. For example, according to the Stanford AI Index Report 2023, the AI landscape has seen significant shifts in publication and collaboration patterns~\cite{AIIndex2023}. Notably, collaborations between educational institutions and industry grew 4.2 times from 2010 to 2021. In 2021, educational institutions and nonprofits led with 32,551 collaborations, followed by industry and educational institutions (12,856), and educational and government institutions (8,913). Furthermore, the number of AI journal publications increased approximately 2.3 times since 2015, showing a 14.8\% rise from 2020 to 2021.

Organizations are also increasingly dependent on ML systems for critical insights and decision-making, highlighting the crucial role of MLOps (Machine Learning Operations) pipelines in achieving reliable and scalable ML implementations. MLOps aims to standardize development, boost productivity, and enable reproducible techniques for rapid experimentation and model training.

\section{Research Problem}

Despite rapid advancements in ML, iterative experimentation workflows face significant computational inefficiencies, particularly in cross-validation scenarios. Leave-One-Subject-Out (LOSO) cross-validation, commonly used to evaluate model generalization across individuals, requires training \textit{N} models for \textit{N} subjects—a process that can take 8-10 hours for datasets with 128 subjects. When researchers iterate on hyperparameters or feature selection strategies, this entire process must be repeated, even when many trained models remain valid under the new configuration.

This research addresses three interconnected challenges in managing LOSO experimentation workflows:

\begin{itemize}

\item \textbf{Redundant Computation Across Experimental Iterations:} 
ML experimentation requires repeatedly modifying specific components 
(hyperparameters, feature subsets, model architectures) while keeping others 
constant. Existing MLOps tools such as DVC~\cite{dvc2022inproceedings} and 
ZenML provide hierarchical pipeline caching, where downstream changes do not 
invalidate upstream artifacts. However, adapting these tools to LOSO workflows 
requires explicit parameterization of the held-out subject and careful cache 
key design, a process that is neither standardized nor validated in existing 
literature. Without a validated framework, researchers face a choice between 
no caching (prohibitive recomputation costs) or manual cache management 
(error-prone and difficult to reproduce).


\item \textbf{Ensuring Cache Validity Across Cross-Validation Folds:}
In LOSO cross-validation, each fold's trained model is valid only for 
evaluating its specific held-out subject. A model trained with subject $S_i$ 
held out produces scientifically invalid predictions when applied to evaluate 
subject $S_j$ held out, not due to data leakage (the train/test split is 
maintained) but due to applying cached results to an incorrect experimental 
condition.

While general-purpose caching systems like DVC can be configured to handle this 
by parameterizing the held-out subject, there is no validated framework 
demonstrating best practices for LOSO cache key design. Without explicit 
guidance, researchers risk either overly conservative cache invalidation 
(recomputing valid results) or incorrect parameterization (potential cache 
validity violations). This research establishes a validated approach where 
held-out subject identifiers are systematically included in fingerprint-based 
cache keys.



\item \textbf{Unified Fingerprinting for Reproducibility and Cache Validity:}
Existing MLOps tools separate experiment tracking (MLflow, W\&B) from 
pipeline caching (DVC, ZenML), requiring researchers to manually coordinate 
logged configurations with cached artifacts~\cite{zaharia2018mlflow,%
wandb2024docs,dvc2022inproceedings}. This separation becomes particularly 
challenging in LOSO scenarios where fold identity must be consistently tracked 
across both systems. This research proposes cryptographic fingerprinting 
(SHA-256 of canonical JSON configurations) as a unified solution: a single 
fingerprint simultaneously ensures cache validity, reproducibility, and 
experiment identification. The approach is validated by demonstrating 
deterministic fingerprint generation across 18 experimental configurations 
and 128 cross-validation folds, with zero fingerprint collisions and 
complete result reproducibility.

\end{itemize}



This thesis addresses these challenges through a fingerprint-based caching 
framework for LOSO cross-validation experiments. The framework uses SHA-256 
hashing of canonical JSON configuration representations to generate deterministic 
cache keys, with held-out subject identifiers included to ensure cache validity 
across folds. The approach is validated through a practical case study: sleep 
stage classification using 128 subjects from the BOAS (Bitbrain Open Access 
Sleep) dataset~\cite{kemp2000sleepEDF,goldberger2000physionet}. Evaluation 
metrics include cache hit rates, computational time savings, storage requirements, 
and reproducibility validation across experimental iterations.


\section{Research Questions}

\textbf{Main Research Question:}

How can fingerprint-based caching reduce computational costs in Leave-One-Subject-Out cross-validation while ensuring reproducibility and preventing data leakage?

\begin{itemize}
\item \textbf{RQ1 (Efficiency):} What time savings and cache hit rates can be achieved through fingerprint-based caching of extracted features and trained models in LOSO cross-validation?
\item \textbf{RQ2 (Reproducibility):} How does SHA-256 configuration fingerprinting ensure consistent and reproducible experiment results across multiple runs?
\item \textbf{RQ3 (Scalability):} How does caching performance scale from pilot studies (10 subjects) to full datasets (128 subjects)?
\item \textbf{RQ4 (Trade-offs):} What are the storage costs versus computational savings of the proposed caching approach?
\end{itemize}



\section{Research Method and Approach}
This research employs a quantitative experimental methodology combining:

\begin{itemize}
\item \textbf{Performance Measurement:} Systematic collection of timing metrics, cache hit rates, and storage utilization across multiple experimental scenarios.
\item \textbf{Comparative Analysis:} Evaluation of caching efficiency across different conditions (cold start vs. cached runs, identical vs. modified configurations).
\item \textbf{Reproducibility Validation:} Verification that identical fingerprints produce identical results across multiple runs.
\end{itemize}

The approach centers on two orthogonal evaluation dimensions:
\begin{itemize}
\item \textbf{WHY to cache:} Optimization goals including latency reduction, throughput improvement, and cumulative time savings across iterative experiments.
\item \textbf{WHERE to cache:} The framework implements hierarchical caching at two levels. Feature-level caching stores 149 extracted features per epoch, computed once per subject (implemented, achieving 227$\times$ speedup in preliminary tests). LOSO model caching will store trained models per cross-validation fold, with the held-out subject identifier encoded in the cache key to prevent data leakage (currently in implementation).
\end{itemize}

\section{Structure of the Thesis}
This thesis is structured as follows:

\begin{itemize}
\item \textbf{Chapter 1: Introduction} presents the motivation, research problem, and main research question for fingerprint-based caching in LOSO cross-validation

\item \textbf{Chapter 2: Background} covers ML research methodology, computational efficiency challenges, and caching technologies.

\item \textbf{Chapter 3: Related Works} reviews existing MLOps tools (MLflow, DVC, ZenML) and sleep stage classification methods, identifying gaps in LOSO-aware caching that this work addresses.

\item \textbf{Chapter 4: Methodology} details the fingerprint-based caching framework, the Sleep-EDF dataset, preprocessing pipeline, extraction of 149 EEG features per epoch, and the LOSO cross-validation experimental design.

\item \textbf{Chapter 5: Results} presents cache performance metrics (hit rates, time savings, storage overhead - currently in inplementation), reproducibility validation, and classification results validating framework correctnes.

\item \textbf{Chapter 6: Conclusion and Future Work} summarizes contributions in fingerprint-based caching for cross-validation experiments and discusses extensions including cross-experiment cache reuse.
\end{itemize}

% ============================================================================
% CHAPTER 2: BACKGROUND (Keep existing content)
% ============================================================================
\chapter{Background}

\section{Domain: Machine Learning Research Methodology}

The landscape of Machine Learning (ML) research and development is 
characterized by rapid iteration and experimentation. According to the 
AI Index Report 2023, AI-related publications have more than doubled since 
2015, reflecting the field's exponential growth~\cite{AIIndex2023}. Data 
scientists and ML engineers continuously explore different model architectures, 
algorithms, hyperparameters, and data preprocessing techniques to achieve 
optimal performance. This iterative process, while essential for innovation, 
introduces significant challenges concerning efficiency and reproducibility. 
A single experiment may require hours or days of computation, and researchers 
often repeat similar computations across multiple runs. The cumulative effect 
translates into substantial costs in GPU hours, cloud computing expenses, 
and researcher productivity~\cite{schubotz2022caching}.

To manage this complexity, standardized process models have emerged. 
CRISP-DM (Cross-Industry Standard Process for Data Mining), developed in 
1996, remains widely adopted~\cite{wirth2000crispdm}. However, its limitations 
for modern ML—particularly lack of emphasis on model monitoring and 
maintenance, led to CRISP-ML(Q), which extends CRISP-DM with quality assurance 
methodology tailored for ML applications~\cite{studer2021crispmlq}. This 
framework aligns with modern MLOps principles, which emphasize automation, 
continuous integration/deployment (CI/CD), and reproducibility throughout 
the ML lifecycle~\cite{kreuzberger2023mlops}.

Tools such as MLflow~\cite{zaharia2018mlflow} and MLXP~\cite{arbel2024mlxp} 
support experiment tracking by recording parameters, metrics, code versions, 
and artifacts for each run, enabling comparison and reproducibility. Despite 
these advances, reproducibility remains a pervasive challenge in ML 
research~\cite{semmelrock2023reproducibility}. Reproducibility ensures consistent 
results from an ML workflow under identical conditions, hinging on the 
``Holy Trinity'': input data, code and parameters, and the execution 
environment~\cite{schubotz2022caching}. Challenges arise from:

\begin{itemize}
    \item \textbf{Lack of Records and Hyperparameter Inconsistency:} 
          Unrecorded decisions or parameter changes hinder replication.
    \item \textbf{Data and Environment Volatility:} Dataset alterations, 
          framework updates, and inconsistent compute environments cause 
          divergent outcomes~\cite{idowu2024empirical}.
    \item \textbf{Inherent Randomness:} Random initializations, data 
          shuffling, and non-deterministic GPU operations make consistent 
          replication challenging~\cite{semmelrock2023reproducibility}.
\end{itemize}

These challenges are particularly acute during iterative experimentation, 
where researchers may run hundreds of configurations before identifying 
promising approaches. Without intelligent management of computational 
artifacts, substantial resources are wasted on redundant calculations. 
This observation motivates the central contribution of this thesis: a 
fingerprint-based caching framework that persists and reuses computational 
artifacts (extracted features and trained models) across iterative modeling 
phases while maintaining reproducibility through cryptographic configuration 
fingerprinting.

\section{Context: Computational Efficiency in Research}

Building on the iterative experimentation challenges described in 
Section 2.1, this section examines the computational efficiency bottlenecks 
that motivate intelligent caching frameworks.

The increasing complexity and scale of Machine Learning models have 
substantially raised computational demands. Training models, tuning 
hyperparameters, and running repeated experiments consume vast compute, 
memory, and time resources, creating bottlenecks in research productivity 
and driving up costs.

Several factors drive this resource consumption:

\begin{itemize}
\item \textbf{Large Datasets:} Massive datasets require intensive I/O and 
      memory usage.
      
\item \textbf{Complex Models:} Both deep neural networks and ensemble 
      methods (e.g., gradient boosting with thousands of trees, large 
      random forests) require substantial computational resources for 
      training and inference.
      
\item \textbf{Iterative Experimentation:} Frequent experiments often repeat 
      computations unnecessarily when modifying a subset of parameters while 
      keeping others constant.
      
\item \textbf{Hyperparameter Optimization:} Searching large hyperparameter 
      spaces is inherently compute-intensive.
      
\item \textbf{Rigorous Validation:} Cross-validation strategies (e.g., 
      k-fold, leave-one-out) multiply training costs by the number of folds.
\end{itemize}

While caching can help, traditional caching systems (e.g., general-purpose 
key-value stores like Redis or Memcached) often fail to meet ML-specific 
needs~\cite{krishna2025cache}. These systems are typically static and focus 
on reducing latency rather than adapting to experiment configuration changes 
or optimizing for ML-specific access patterns. Such inefficiencies can slow 
research progress and increase costs.

This thesis addresses these inefficiencies through a fingerprint-based 
caching framework specifically designed for iterative ML workflows. By 
encoding complete experiment configurations (including cross-validation 
fold identity) into cryptographic cache keys, the framework enables 
automatic result reuse while maintaining scientific validity and 
reproducibility.

\section{Technology: Configuration Fingerprinting and Caching}

\begin{center}
\colorbox{yellow!20}{%
\begin{minipage}{0.9\textwidth}
\textbf{Draft Note:} Some statistics and system-specific claims in this section 
are from papers I reviewed, and written down notes but haven't yet added to the bibliography. These 
need proper bib entries and source verification before final submission. 
Marked with \textcolor{red}{[TODO: Add bib entry]}.
\end{minipage}
}
\end{center}

Efficient ML experimentation depends on configuration management and intelligent 
result reuse. This section reviews technologies directly relevant to the 
fingerprint-based LOSO caching approach developed in this thesis.

\subsection{Configuration Fingerprinting}

Configuration fingerprinting uses cryptographic hash functions (MD5, SHA-256) 
to create unique identifiers for experiment configurations. By hashing a 
canonical representation of all parameters, these fingerprints enable 
deterministic identification of identical configurations and automatic 
cache invalidation when parameters change.

Modern MLOps tools employ this approach:
\begin{itemize}
\item \textbf{DVC} logs pipeline signatures combining dependency contents, 
      parameters, and commands~\cite{dvc2022inproceedings}
\item \textbf{MLflow} tracks experiment metadata and enables visual 
      comparison~\cite{zaharia2018mlflow}
\item \textbf{MLXP} uses Hydra-based configuration logging with Git 
      hashes~\cite{arbel2024mlxp}
\end{itemize}

\subsection{Hash-Based Caching in MLOps}

MLOps pipeline caching implements content-addressable storage: when pipeline 
inputs and code match previous runs (verified through hash comparison), 
cached outputs are reused~\cite{dvc2022inproceedings}. DVC's run cache, 
for example, combines dependency hashes with parameter configurations to 
create cache keys. When a match is found, execution is skipped and previous 
results are restored.

This approach provides reproducibility: identical configurations produce 
identical cache keys, ensuring cached results correspond to recorded 
parameters~\cite{schubotz2022caching}.

\subsection{Related Caching Approaches}

\begin{center}
\colorbox{red!20}{%
\begin{minipage}{0.9\textwidth}
\textcolor{red}{\textbf{TODO:}} Add bib entries for: vLLM/PagedAttention paper 
(60-80\% memory reduction claim), GATI learned caching paper (7.69$\times$ speedup), 
XRootD/XCache infrastructure reports (10.8 PB statistic). Papers are in my 
download folder but not yet in references.bib.
\end{minipage}
}
\end{center}

While this thesis focuses on hash-based caching for ML experiments, other 
caching paradigms exist in adjacent domains:

\textbf{Large Language Model (LLM) Inference Caching} addresses inference 
latency by caching attention key-value tensors during autoregressive generation. 
Systems like vLLM's PagedAttention reportedly reduce memory waste from 60--80\% 
to under 4\% and enable up to 24$\times$ higher throughput \textcolor{red}{[TODO: Add bib entry]}. 
Prefix caching can achieve up to 85\% latency reduction and 90\% cost savings 
for repeated system prompts \textcolor{red}{[TODO: Add bib entry]}. These techniques 
are specific to LLM inference and not directly applicable to experiment result caching.

\textbf{Learned Caching for Deep Networks} employs neural networks to predict 
hidden layer outputs, enabling layer skipping during inference. GATI, for 
example, achieves up to 7.69$\times$ latency reduction at 97\% accuracy 
\textcolor{red}{[TODO: Add bib entry]}. This represents a different caching 
philosophy—using ML to make cache decisions—but applies only to inference scenarios.

\textbf{High-Energy Physics (HEP) Data Caching} operates at petabyte scale 
for raw data files. Systems like XRootD/XCache achieve cache hit rates exceeding 
85\% by request count. The Southern California Cache infrastructure reportedly 
saved 10.8 PB of network traffic in a single year \textcolor{red}{[TODO: Add bib entry]}. 
However, these systems are optimized for data distribution, not experiment 
configuration management.

In contrast to these approaches, this thesis implements deterministic, 
hash-based caching specifically designed for cross-validation experiments, 
where cache validity depends on both configuration parameters and fold identity.

\subsection{Reproducibility Considerations}

Despite fingerprinting, perfect reproducibility remains challenging. Sources 
of non-determinism include random seed variation, software version differences, 
hardware-specific floating-point operations, and non-deterministic GPU 
algorithms~\cite{semmelrock2023reproducibility}.

\begin{center}
\colorbox{red!20}{%
\begin{minipage}{0.9\textwidth}
\textcolor{red}{\textbf{TODO:}} Verify 0.2--0.5\% variance statistic. Check if 
this is in Semmelrock 2023 or requires separate citation (possibly Pham et al. 
2020 on variance in deep learning systems).
\end{minipage}
}
\end{center}

Empirical studies suggest that even with controlled random seeds, small 
variance (approximately 0.2--0.5\%) may persist between runs due to hardware-level 
floating-point non-determinism \textcolor{red}{[TODO: Add bib entry]}. Configuration 
fingerprinting addresses controllable factors (seeds, software versions, data 
splits) while accepting that hardware-induced variance cannot be eliminated 
through software means. For caching purposes, cached results are considered 
valid if configuration fingerprints match, with the understanding that minor 
numerical differences may occur due to non-deterministic hardware operations.
% ============================================================================
% CHAPTER 3: RELATED WORKS (Keep existing content)
% ============================================================================
\chapter{Related Works}

This chapter reviews existing solutions for ML experiment caching and identifies 
the specific gap this thesis addresses: the lack of validated best practices 
for caching in Leave-One-Subject-Out (LOSO) cross-validation scenarios.

\section{Academic Research on ML Caching}

Recent research explores adaptive caching strategies for ML workloads. Krishna 
(2025) surveys ML-aided caching techniques including LSTM-based cache usage 
prediction and reinforcement learning for adaptive cache policies~\cite{krishna2025cache}. 
These approaches focus on \textit{learned} caching decisions using predictive 
models.

In the context of reproducibility, Schubotz et al. (2022) demonstrate how 
caching can accelerate data science experiments while maintaining FAIR 
principles~\cite{schubotz2022caching}. Their work emphasizes content-addressable 
storage for automatic result reuse, similar to the fingerprinting approach 
used in this thesis.

However, existing academic work on ML caching does not specifically address 
the requirements of cross-validation experiments, where cache validity depends 
on explicit encoding of data splits. The challenge of fold-aware cache key 
generation remains unaddressed in the literature.

\section{Commercial MLOps Tools}

Current MLOps platforms provide general experiment tracking and pipeline 
caching but lack native support for cross-validation-specific scenarios:

\begin{itemize}
\item \textbf{MLflow:} Provides experiment tracking, model registry, and 
      basic result caching~\cite{zaharia2018mlflow}. Pipeline caching requires 
      the MLflow Recipes module, which caches individual steps but does not 
      natively encode fold assignments in cache keys.

\item \textbf{DVC:} Uses MD5 hashes of pipeline inputs and code to enable 
      pipeline stage reuse~\cite{dvc2022inproceedings}. While DVC can cache 
      cross-validation results, it requires manual configuration of the pipeline 
      definition to parameterize fold indices—a process not documented in 
      standard DVC workflows.

\item \textbf{ZenML:} Implements hash-based step caching with automatic 
      invalidation when inputs or code change. However, ZenML's caching 
      abstractions do not provide explicit support for encoding held-out 
      subject identifiers, requiring users to implement custom materializers.

\item \textbf{Neptune AI and Weights \& Biases:} Excel at experiment logging, 
      metric visualization, and collaborative features~\cite{neptune2024docs,wandb2024docs}, 
      but do not provide computational caching mechanisms—only metadata tracking.
\end{itemize}

\section{Research Gap Analysis}

Table~\ref{tab:related_works_comparison} compares caching capabilities of 
existing tools relevant to this thesis.


\begin{center}
\colorbox{red!20}{%
\begin{minipage}{0.9\textwidth}
\textcolor{red}{\textbf{TODO:}} Confirm on latest tool versions + format table new
\end{minipage}
}
\end{center}


\begin{table}[ht]
\centering
\begin{tabular}{|p{2.5cm}|p{2.5cm}|p{2.5cm}|p{5cm}|}
\hline
\textbf{Tool/Framework} & \textbf{Caching Strategy} & \textbf{Fingerprinting Method} & \textbf{LOSO Support} \\
\hline
MLflow & Step-level (optional) & Parameter hashing & Manual fold parameterization \\
\hline
DVC & Pipeline stage & MD5 (code + inputs) & Manual pipeline params \\
\hline
ZenML & Step-level & SHA-256 (inputs + code + env) & Custom materializers \\
\hline
Neptune AI / W\&B & None (tracking only) & N/A & N/A \\
\hline
\textbf{This Thesis} & \textbf{Hierarchical (features + models)} & \textbf{SHA-256 with subject ID} & \textbf{Native fold-aware keys} \\
\hline
\end{tabular}
\caption{Comparison of caching capabilities in MLOps tools. ``LOSO Support'' 
indicates whether the tool natively encodes fold assignments in cache keys 
or requires manual configuration.}
\label{tab:related_works_comparison}
\end{table}

The core gap is the lack of validated, documented patterns for 
implementing fold-aware caching in cross-validation workflows. While tools 
like DVC and ZenML provide general pipeline caching infrastructure, adapting 
them for LOSO requires custom logic that:
\begin{enumerate}
\item Encodes the held-out subject identifier in cache keys to prevent 
      cross-fold contamination
\item Manages hierarchical dependencies between preprocessing, feature 
      extraction, and model training stages
\item Validates cache hits against both parameter configurations and 
      fold assignments
\end{enumerate}

These requirements are not addressed in standard documentation, and incorrect 
implementation risks subtle bugs where models trained on Subject A are 
erroneously reused for evaluating Subject B.

This thesis provides a reference implementation demonstrating these patterns 
in a real LOSO scenario (128-fold cross-validation on Sleep-EDF), with 
explicit measurement of cache performance (hit rates, time savings) and 
validation of reproducibility. The contribution is not a new MLOps platform, 
but a validated approach for extending existing fingerprinting 
techniques to cross-validation contexts.

% ============================================================================
% CHAPTER 4: METHODOLOGY (SIGNIFICANTLY EXPANDED)
% ============================================================================
\chapter{Methodology}

This chapter presents the complete methodology for implementing and evaluating the intelligent caching framework. Section~\ref{sec:theoretical} establishes the theoretical foundations, Section~\ref{sec:data} describes the dataset, Sections~\ref{sec:preprocessing}--\ref{sec:features} detail the signal processing and feature extraction pipeline, Section~\ref{sec:caching} explains the caching implementation, Section~\ref{sec:experiment_setup} outlines the experimental design, and Section~\ref{sec:evaluation_plan} defines the evaluation metrics.

\section{Theoretical Framework}
\label{sec:theoretical}

The proposed framework is grounded in three key concepts: configuration fingerprinting, a distinction between persistent storage and intelligent caching, and two orthogonal evaluation dimensions. This section presents both the currently implemented components and the planned extensions that form the core contribution of this thesis.

\subsection{Configuration Fingerprinting}

Configuration fingerprinting provides the foundation for cache key generation and reproducibility tracking. Each experiment configuration is uniquely identified by a cryptographic hash computed as:

\begin{equation}
\text{fingerprint} = \text{SHA256}(\text{canonical\_json}(\text{config}))[:32]
\end{equation}

The fingerprint uses 32 hexadecimal characters (128 bits), providing strong collision resistance suitable for large-scale experiment tracking. The configuration dictionary includes:

\begin{itemize}
    \item \textbf{Random seed:} Ensures deterministic behavior (e.g., \texttt{42})
    \item \textbf{Code version:} Git commit or semantic version (e.g., \texttt{v1.0}, \texttt{git:abc123})
    \item \textbf{Model configuration:} Algorithm name and hyperparameters (e.g., \texttt{xgboost}, \texttt{max\_depth: 6})
    \item \textbf{Feature configuration:} Base feature count, correlation threshold, and top-K selection
    \item \textbf{Held-out subject ID:} Critical for LOSO---prevents cross-fold contamination (e.g., \texttt{SC4001})
\end{itemize}

An example fingerprint: \texttt{a3f2b8c1d4e5f6a7b8c9d0e1f2a3b4c5}.

The fingerprint generation module is fully implemented, providing the \texttt{LOSOFingerprint} class with methods for generating fingerprints from configuration objects or dictionaries. The \texttt{generate\_cache\_key()} function provides a convenient interface that combines the fingerprint with the held-out subject ID (e.g., \texttt{a3f2b8c1...\_SC4001}) for direct use as cache file names.

\subsection{Configuration-Independent vs. Configuration-Dependent Caching}
\label{sec:storage_vs_caching}

A fundamental distinction in this framework is between \textit{configuration-independent caching} (artifacts that depend only on the dataset) and \textit{configuration-dependent caching} (artifacts that require fingerprint-based invalidation when experimental parameters change). While both avoid redundant computation, they serve different purposes:

\textbf{Configuration-Independent Cache} refers to artifacts that are computed once per dataset version and remain valid across all experimental variations. These artifacts do not depend on model hyperparameters, feature selection criteria, or cross-validation strategies. In this framework, the extracted feature matrices constitute configuration-independent caching---they are computed once per subject and reused across all experiments.

\textbf{Configuration-Dependent Cache} refers to artifacts whose validity depends on experimental configuration. These artifacts require fingerprint-based invalidation when parameters change. The LOSO model cache, which stores trained models keyed by configuration fingerprint and held-out subject, exemplifies configuration-dependent caching.

Table~\ref{tab:storage_hierarchy} summarizes this distinction:

\begin{center}
\colorbox{yellow!20}{%
\begin{minipage}{0.9\textwidth}
\textbf{Working Draft Note:} 
currently in last development Status and finding bugs + fixing table layout
\end{minipage}
}
\end{center}
\begin{table}[ht]
\centering
\begin{tabular}{|l|l|l|l|l|}
\hline
\textbf{Layer} & \textbf{Artifact} & \textbf{Size/Subject} & \textbf{Type} & \textbf{Status} \\
\hline
Layer 1 & 195 features (NPZ) & $\sim$1.1 MB & Persistent Storage & \textbf{Implemented} \\
\hline
Layer 2 & Trained LOSO models & $\sim$10--50 MB & Fingerprint Cache & \textbf{Implemented \& currently verfied} \\
\hline
\end{tabular}
\caption{Storage hierarchy distinguishing persistent storage from intelligent 
caching. Both layers are operational; verification testing demonstrated 
100\% cache hit rate and 4.5$\times$ speedup for identical experiment re-runs.}
\label{tab:storage_hierarchy}
\end{table}

\subsubsection{Layer 1: Feature Extraction as Persistent Storage (Implemented)}

The feature extraction layer computes \textbf{all 195 features} for each subject (representing the full 8-channel configuration: 6 EEG + EOG + EMG) and stores them in compressed NPZ format. This is classified as persistent storage because:

\begin{itemize}
    \item Features depend only on the raw EEG data and fixed preprocessing parameters
    \item Changing model hyperparameters or feature selection thresholds does not invalidate the stored features
    \item The full feature set is extracted once; downstream selection operates on this cached superset
\end{itemize}

\textbf{Feature Count Clarification.} The cache stores 195 features (23 per-channel features $\times$ 8 channels + 11 global features). For experiments using only the 6 EEG channels, features are filtered to 149 at query time ($<$1 second) via the \texttt{select\_channel\_features()} function, avoiding cache invalidation when switching between channel configurations.

This layer is fully implemented and provides measured speedups of approximately 224$\times$ for feature loading compared to re-extraction.

\subsubsection{Layer 2: LOSO Model Caching (Primary Contribution)}
\label{sec:loso_caching}

The LOSO model cache represents the \textbf{primary contribution} of this thesis. 
The complete system—including fingerprint generation, cache storage, retrieval logic, and automatic invalidation, is in implementation right now and validated through comprehensive testing (see Section~\ref{sec:verification}).



\begin{center}
\colorbox{yellow!20}{%
\begin{minipage}{0.9\textwidth}
\textbf{Working Draft Note:} 
As of today 13.01.2026 this is the current implementation status
\end{minipage}
}
\end{center}

\textbf{Implementation Status:} The caching system achieved 100\% cache hit rate 
in verification tests, with measured speedup of 4.5$\times$ for identical re-runs. 
Testing across 24 training configurations (3 subjects $\times$ 8 experimental 
combinations) demonstrated zero cache corruption errors and correct fingerprint-based 
invalidation when parameters changed.

\textbf{Why LOSO-Level Caching, Not Just Model Caching?}

A critical design decision is caching at the \textit{LOSO fold level} rather than caching a single global model. This distinction is essential:

\begin{enumerate}
    \item \textbf{Data Leakage Prevention:} Each LOSO fold has a unique fingerprint that includes the \texttt{held\_out\_subject} field. A model trained with Subject~A held out \textit{cannot} be retrieved when Subject~B is the test subject---the fingerprints differ (as demonstrated by the implemented \texttt{LOSOFingerprint} class), forcing a cache miss and retraining.
    
    \item \textbf{Scientific Validity:} In LOSO cross-validation, each fold produces a \textit{different} model because the training data differs. Caching a single ``global model'' would be scientifically invalid---there is no single model, but rather $N$ distinct models (one per subject held out).
    
    \item \textbf{Incremental Experimentation:} When adding new subjects or modifying experimental parameters, only affected folds require retraining. Cached models for unchanged configurations remain valid and are reused automatically.
\end{enumerate}

\textbf{Why Model-Level Caching, Not Result-Level Caching?}

Within each LOSO fold, the framework caches the \textit{trained model} rather than just the evaluation metrics:

\begin{itemize}
    \item \textbf{Post-hoc Analysis:} Cached models enable feature importance analysis, prediction on new data, or computation of additional metrics without retraining.
    \item \textbf{Reproducibility Verification:} Identical fingerprints guarantee identical model states, enabling verification that predictions (not just aggregate metrics) are reproducible.
    \item \textbf{Debugging Capability:} When unexpected results occur, the exact model state can be inspected rather than requiring reproduction of the entire training process.
\end{itemize}

The trade-off is increased storage (models are larger than metric summaries), but this is acceptable given the scientific benefits and modern storage capacities.

\subsection{Two Orthogonal Evaluation Dimensions}

The framework evaluation follows two orthogonal dimensions that structure the experimental analysis:

\begin{itemize}
    \item \textbf{WHY dimension (Optimization Goals):}
    \begin{itemize}
        \item Latency: Time to retrieve cached vs. computed results
        \item Throughput: Subjects processed per unit time
        \item Time Saved: Cumulative computation time avoided across experimental iterations
    \end{itemize}
    \item \textbf{WHERE dimension (Cached Artifact Levels):}
    \begin{itemize}
        \item Layer 1: Extracted feature matrices (persistent storage---implemented)
        \item Layer 2: Trained LOSO fold models (intelligent cache---in progress)
        \item Results: Final metrics (always recomputed from cached models for verification)
    \end{itemize}
\end{itemize}

This two-dimensional framework enables systematic evaluation of caching benefits across different experimental scenarios, from simple re-runs (maximum cache hits) to parameter sweeps (partial cache hits) to completely new configurations (cache misses establishing new baselines).

\section{Data: BOAS Dataset}

\label{sec:data}

\subsection{Dataset Overview}

This research uses the Bitbrain Open Access Sleep (BOAS) dataset, a publicly available polysomnographic EEG database. Table~\ref{tab:dataset_overview} summarizes the dataset characteristics.

\begin{table}[ht]
\centering
\begin{tabular}{|l|l|}
\hline
\textbf{Parameter} & \textbf{Value} \\
\hline
Dataset & Bitbrain Open Access Sleep (BOAS) \\
Total Subjects & 128 healthy adults \\
Epochs per Subject & $\sim$940 (varies by recording length) \\
Total Epochs & $\sim$120,000 \\
EEG Channels & 6 (F3, F4, C3, C4, O1, O2) following 10-20 system \\
Sampling Rate & 256 Hz (downsampled to 128 Hz after preprocessing) \\
Epoch Duration & 30 seconds \\
Samples per Epoch & 3,840 (at 128 Hz) \\
Sleep Stages & 5 (Wake, N1, N2, N3, REM) per AASM criteria \\
\hline
\end{tabular}
\caption{BOAS dataset characteristics.}
\label{tab:dataset_overview}
\end{table}

\subsection{Dataset Selection Justification}

The BOAS dataset was selected over Sleep-EDF for several reasons:
\begin{itemize}
\item \textbf{Public Availability:} Freely accessible via OpenNeuro, enabling reproducibility.
\item \textbf{Larger Subject Pool:} 128 subjects with consistent recording quality.
\item \textbf{Complete Channel Coverage:} All required EEG channels available across all subjects.
\item \textbf{Modern Annotations:} Sleep stages annotated according to AASM criteria.
\end{itemize}

\subsection{Expected Class Distribution}

Sleep stage classification exhibits inherent class imbalance, with N1 (light sleep) typically underrepresented. Based on literature and preliminary analysis, the expected distribution is approximately: Wake (15--20\%), N1 (5--10\%), N2 (45--50\%), N3 (15--20\%), REM (15--20\%). This imbalance will be addressed through macro-averaged metrics in evaluation.

\section{Data Validation and Integrity Verification}
\label{sec:data_validation}

Ensuring data quality and integrity is critical for reproducible sleep staging research. Prior to feature extraction, each polysomnographic recording undergoes a comprehensive 12-point validation protocol to detect common data issues that could compromise classification accuracy or introduce systematic biases.

\subsection{Motivation}

Sleep EEG datasets frequently contain artifacts from various sources: electrode disconnections during subject movement, signal clipping from amplifier saturation, flat channels from electrode detachment, or misalignment between signal data and annotation files. Without systematic validation, such issues may go undetected, leading to corrupted features that degrade classifier performance or, worse, produce misleadingly optimistic results when problematic epochs are inadvertently excluded from evaluation.

The BOAS dataset explicitly marks disconnection events (stage code 8) and artifacts (stage code -2) in the annotation files. However, additional integrity checks are necessary to verify temporal alignment between annotations and signal data, detect hardware failures not captured in annotations, and ensure data completeness.

\subsection{Validation Protocol}

The implemented validation protocol comprises 12 automated checks organized into four categories.

\textbf{Annotation Integrity (5 checks):}
\begin{itemize}
\item \textbf{Sufficient epochs:} $n_{\text{valid}} \geq 100$ ensures adequate data for statistical analysis
\item \textbf{All stages present:} $\{\text{W, N1, N2, N3, REM}\} \subseteq \text{recording}$ detects incomplete recordings or labeling errors
\item \textbf{Sufficient duration:} $t \geq 7200\text{s}$ (2 hours) excludes truncated or partial recordings
\item \textbf{Acceptable filter rate:} $\text{filtered}/\text{total} < 30\%$ indicates systematic problems when exceeded
\item \textbf{Temporal continuity:} $\max(\Delta t) \leq 1.5 \times \text{epoch\_duration}$ detects unexpected gaps in annotation timeline
\end{itemize}

\textbf{Structural Consistency (2 checks):}
\begin{itemize}
\item \textbf{Correct channels:} $n_{\text{channels}} \geq 6$ verifies required EEG montage is present
\item \textbf{Correct sampling rate:} $|f_s - f_{\text{target}}| < 1\text{ Hz}$ ensures consistent temporal resolution
\end{itemize}

\textbf{EDF-Annotation Alignment (2 checks):}
\begin{itemize}
\item \textbf{Annotations within EDF:} $\max(\text{onset}) + \text{epoch\_dur} \leq \text{EDF\_duration}$ prevents out-of-bounds data access
\item \textbf{Sample bounds valid:} $\max(\text{endsample}) \leq n_{\text{samples}}$ verifies sample indices reference existing data
\end{itemize}

\textbf{Signal Quality (3 checks):}
\begin{itemize}
\item \textbf{No invalid values:} $\nexists$ NaN $\vee$ $\pm\infty$ detects corrupted or missing samples
\item \textbf{No flat channels:} $\sigma_{\text{channel}} > 10^{-10}$ identifies dead electrodes or recording failures
\item \textbf{Reasonable amplitude:} $\max(|x|) < 10^6$ detects extreme values indicating unit errors or clipping
\end{itemize}

\subsection{Implementation}

Signal quality checks are performed on representative segments (first and last 30 seconds) rather than the complete recording to minimize computational overhead while still detecting common failure modes such as electrode detachment at recording start or end. For each validation failure, a descriptive error message is logged to facilitate debugging and enable transparent reporting of data exclusions.

\subsection{Validation Results: Preliminary Subset}

Table~\ref{tab:validation_results} shows validation results for the first three 
subjects, representing a pilot verification of the data quality protocol. 

\begin{table}[ht]
\centering
\begin{tabular}{|l|c|c|c|l|}
\hline
\textbf{Subject} & \textbf{Valid Epochs} & \textbf{Duration (h)} & \textbf{Filtered} & \textbf{Issues} \\
\hline
sub-1 & 914/915 & 7.63 & 0.1\% & 1 disconnection at $t=3690$s \\
\hline
sub-2 & 913/913 & 7.61 & 0.0\% & None \\
\hline
sub-3 & 997/997 & 8.31 & 0.0\% & None \\
\hline
\end{tabular}
\caption{Validation results for pilot dataset (3/128 subjects). Full dataset 
validation is ongoing; preliminary results confirm high data quality suitable 
for ML analysis.}
\end{table}

\textbf{Status:} As of this writing, comprehensive validation of all 128 subjects 
is in progress. The pilot results demonstrate that the BOAS dataset exhibits 
excellent signal integrity, with minimal artifact contamination ($<$0.1\% epoch 
exclusion rate) and high inter-rater agreement (83.5\% human-AI concordance). 
Complete validation results will be included in the final thesis submission.

All subjects passed the critical signal integrity checks (no NaN values, no flat channels, reasonable amplitudes), confirming that the BOAS dataset provides high-quality polysomnographic recordings suitable for machine learning analysis. The single disconnection event in subject 1 was correctly annotated in the original dataset and automatically excluded during preprocessing.


\subsection{Human vs.\ AI Label Agreement}

An additional quality metric available in the BOAS dataset is the agreement between human expert annotations (\verb|stage_hum|) and automated AI staging (\verb|stage_ai|). For subject 1, the inter-rater agreement was 83.5\%, which aligns with reported inter-scorer reliability in clinical polysomnography (typically 80--90\%). This provides a useful reference point for evaluating classifier performance: achieving accuracy significantly below human-AI agreement may indicate model issues, while exceeding it suggests the classifier captures patterns consistent with  the consensus human annotation standard.

\subsection{Importance for Reproducibility}

By systematically validating data integrity before analysis, this pipeline ensures that:

\begin{enumerate}
\item \textbf{No silent failures:} Data integrity issues are detected and 
      logged explicitly, preventing silent propagation of corrupted features through the pipeline
\item \textbf{Transparent exclusions:} Any filtered epochs or subjects are documented with specific reasons
\item \textbf{Consistent preprocessing:} All subjects undergo identical validation criteria
\item \textbf{Reproducible results:} The validation protocol can be applied to new datasets or by other researchers
\end{enumerate}

This approach aligns with best practices for transparent reporting in machine learning research and addresses concerns about hidden data preprocessing decisions that can affect reported classification accuracy.


\section{Preprocessing Pipeline}
\label{sec:preprocessing}

The preprocessing pipeline transforms raw EDF recordings into clean, epoched signals suitable for feature extraction.

\subsection{Signal Filtering}

The following signal conditioning is applied sequentially to each EEG channel:

\begin{enumerate}
\item \textbf{Bandpass Filter:} A 5th-order Butterworth filter with cutoff frequencies of 0.5--40 Hz removes DC drift and high-frequency noise while preserving sleep-relevant frequency bands.
\item \textbf{Notch Filter:} A notch filter at 50 Hz (Q=30) removes power line interference.
\item \textbf{Downsampling:} 256 Hz $\rightarrow$ 128 Hz using decimation with built-in anti-aliasing low-pass filter (prevents aliasing artifacts).
\end{enumerate}

All preprocessing parameters are included in the cache fingerprint, ensuring automatic invalidation when parameters change.

\subsection{Epoch Creation}

Continuous signals are segmented into 30-second epochs aligned with hypnogram annotations. Each epoch contains 3,840 samples (30 seconds $\times$ 128 Hz) across 6 channels, resulting in an array of shape $(n_\text{epochs}, 3840, 6)$ per subject.

\subsection{Quality Control}

Epochs are filtered according to the comprehensive validation criteria established in Section~\ref{sec:data_validation}. In summary, epochs must satisfy:
\begin{itemize}
\item Amplitude within $\pm$500 $\mu$V (reject artifacts)
\item No NaN or infinite values
\item Valid sleep stage annotation (exclude disconnection events, artifacts, and unknown stages)
\end{itemize}

The full 12-point validation protocol ensures data integrity before feature extraction begins.

\section{Feature Extraction: The 149-Feature Set}
\label{sec:features}

This research employs a comprehensive extraction of 149 EEG features per epoch strategy designed for robust cross-subject generalization. Features are computed for each 30-second epoch and organized into four domains, as detailed in Table~\ref{tab:feature_breakdown_149}.

\begin{table}[ht]
\centering
\begin{tabular}{|p{3.5cm}|p{1.5cm}|p{7cm}|}
\hline
\textbf{Domain} & \textbf{Count} & \textbf{Features} \\
\hline
Time-Domain & 60 & Mean, std, variance, min, max, peak-to-peak, RMS, skewness, kurtosis, zero-crossing rate (10 $\times$ 6 channels) \\
\hline
Frequency-Domain & 54 & Band powers (delta 0.5--4 Hz, theta 4--8 Hz, alpha 8--13 Hz, sigma 12--15 Hz, beta 13--30 Hz, gamma 30--40 Hz), spectral entropy, peak frequency, median frequency (9 $\times$ 6 channels) \\
\hline
Signal Complexity & 24 & Hjorth mobility, Hjorth complexity, Hurst exponent, Detrended Fluctuation Analysis (4 $\times$ 6 channels) \\
\hline
Inter-Channel & 11 & Coherence pairs (6), Phase Locking Value pairs (3), global entropy, global complexity \\
\hline
\textbf{Total} & \textbf{149} & \textbf{Comprehensive EEG-based feature set} \\
\hline
\end{tabular}
\caption{Complete 149-feature breakdown for sleep stage classification.}
\label{tab:feature_breakdown_149}
\end{table}

\subsection{Feature Engineering Strategy and Caching Prioritization}

The feature extraction pipeline implements a comprehensive 149-feature representation spanning time-domain statistics, frequency-domain band powers, complexity measures, and cross-channel connectivity metrics. This feature set was deliberately designed to balance discriminative power for sleep stage classification with computational tractability for caching optimization.

While additional features such as sample entropy, permutation entropy, band power ratios, or spindle density detection could marginally enhance classification performance, the current architecture intentionally prioritizes caching efficiency as the central research contribution. By establishing a fixed, comprehensive feature space, the system enables complete cache validity across experimental variations in feature selection, model hyperparameters, and cross-validation strategies---changes that would otherwise trigger costly recomputation.

This design philosophy reflects a strategic trade-off: rather than pursuing incremental gains in classification accuracy through feature proliferation, the framework optimizes for reproducibility, computational efficiency, and scalability. The result is a system where subsequent experiments achieve near-instantaneous feature retrieval, reducing pipeline latency from hours to seconds, minimizing computational resource consumption, and ensuring deterministic reproducibility across research iterations. Future extensions may incorporate additional feature domains; however, the current implementation demonstrates that intelligent caching of a well-designed feature set delivers substantial practical benefits that outweigh marginal accuracy improvements from feature expansion.

\subsection{Feature Selection Strategy}

After extracting all 149 features, optional feature selection is applied using a two-stage pipeline:

\textbf{Stage 1: Correlation-Based Filtering}
\begin{itemize}
\item Removes highly correlated features using Pearson correlation
\item Threshold options: 0.75, 0.90, or None (no filtering)
\item Prevents multicollinearity in downstream models
\end{itemize}

\textbf{Stage 2: ANOVA F-statistic Selection}
\begin{itemize}
\item Ranks features by class-discriminative power using F-statistic (\texttt{f\_classif})
\item Top-K options: 30, 50, or None (retain all 149 features)
\item ANOVA selected over Mutual Information due to approximately 76$\times$ faster computation with minimal accuracy difference (0.2\% in preliminary benchmarks: ANOVA 80.1\% vs.\ MI 80.3\%)\footnote{Timing benchmarks: ANOVA 0.078s vs.\ MI 5.956s for top-50 feature selection on the 5-subject pilot dataset (benchmark conducted December 2025).}
\end{itemize}

\textbf{Implementation Note:} Feature selection is performed \textbf{globally} (not per-fold) to enable consistent cache fingerprinting. This design choice prioritizes caching efficiency while accepting minimal impact on cross-validation purity. The Stage 2 cache stores all 149 features, and selection parameters only affect the final model training step.

\subsection{Feature Extraction Optimization}

To enable efficient processing of the BOAS dataset (128 subjects, $\sim$1000 epochs each, $\sim$128,000 total epochs), several optimization strategies were evaluated and implemented.

\subsubsection{Implemented Optimizations (No Quality Loss)}

\textbf{Multiprocess Parallelization:} Feature extraction was parallelized across CPU cores using Python's joblib library with process-based workers. Thread-based parallelization was initially attempted but provided no speedup due to Python's Global Interpreter Lock (GIL) blocking concurrent execution of CPU-bound NumPy/SciPy operations. Process-based parallelization achieved near-linear scaling, reducing per-epoch extraction time from $\sim$4.7s to $\sim$1.1s on an 8-core system (4.3$\times$ speedup).

\textbf{Optimized Complexity Measures:} The computationally expensive complexity features (Detrended Fluctuation Analysis, Hurst exponent, Hjorth parameters) were accelerated using the antropy library, which provides Numba JIT-compiled implementations of identical peer-reviewed algorithms. This optimization maintained mathematical equivalence while reducing complexity feature computation time by approximately 2--3$\times$.

\textbf{Feature Caching Strategy:} A comprehensive caching system was implemented to compute features once for all 6 EEG channels and store them in compressed NumPy archives. Subsequent experiments can load cached features and filter to the required channel subset at runtime, eliminating redundant computation when comparing channel configurations.

\subsubsection{Evaluated But Not Implemented (Would Reduce Quality)}

\textbf{Approximate Hurst Exponent:} Replacing the DFA-based Hurst exponent with faster auto\-correlation-based approximations was considered but rejected. While offering 5--10$\times$ speedup, this approach produces systematically different values (10--15\% deviation) and would compromise comparability with existing sleep staging literature that relies on standard DFA.

\textbf{Reduced DFA Scale Resolution:} Decreasing the number of window scales in DFA from 10 to 5 was evaluated. Although this would halve DFA computation time, it introduces 2--3\% variance in the resulting exponents, potentially affecting the discrimination of sleep stages with similar spectral profiles (particularly N2/N3 differentiation).

\textbf{Sample Entropy Approximation:} Fast approximate entropy measures were considered but not implemented, as sample entropy's $O(N^2)$ complexity was already addressed through parallelization, and approximate methods would alter the feature's discriminative properties for sleep stage classification.

\subsubsection{Performance Summary}

Table~\ref{tab:optimization_summary} summarizes the optimization decisions and their impact.

\begin{table}[ht]
\centering
\begin{tabular}{|l|c|c|c|}
\hline
\textbf{Optimization} & \textbf{Status} & \textbf{Speedup} & \textbf{Quality Impact} \\
\hline
Process parallelization & Implemented & 4.3$\times$ & None \\
\hline
Antropy complexity & Implemented & 2--3$\times$ & None \\
\hline
Feature caching & Implemented & N/A* & None \\
\hline
Approximate Hurst & Rejected & 5--10$\times$ & $-$10--15\% accuracy \\
\hline
Reduced DFA scales & Rejected & 1.5$\times$ & $-$2--3\% variance \\
\hline
\end{tabular}
\caption{Feature extraction optimization decisions. *Caching eliminates recomputation entirely for repeated experiments.}
\label{tab:optimization_summary}
\end{table}

The combined optimizations (process-based parallelization on 8 cores + Numba-accelerated complexity measures via antropy) reduced total feature extraction time from approximately 150 hours (sequential, unoptimized) to approximately 53 minutes for the complete 128-subject dataset. With the caching system enabled, subsequent runs retrieve cached features near-instantly ($<$1 minute), representing a 227$\times$ speedup over the cold start.



\section{Caching Framework Implementation}
\label{sec:caching}

\subsection{Fingerprint Generation}

Cache keys are generated using deterministic JSON serialization followed by SHA-256 hashing:

\begin{verbatim}
import hashlib
import json

def generate_fingerprint(config: dict) -> str:
    """Generate SHA-256 fingerprint from configuration."""
    canonical = json.dumps(config, sort_keys=True, default=str)
    return hashlib.sha256(canonical.encode()).hexdigest()[:32]
\end{verbatim}

\subsection{Cache Storage Structure}

The cache uses a directory structure organized by fingerprint:

\begin{verbatim}
cache/
+-- features/ # Extracted features (all 195, config-independent)
| +-- subject_sub-1_full.npz  # NO fingerprint - always valid
| +-- subject_sub-2_full.npz
| +-- ...
+-- loso_model_cache/ # Trained models (config-dependent)
  +-- {fingerprint}_sub-1.joblib  # WITH fingerprint
  +-- {fingerprint}_sub-2.joblib
  +-- cache_registry.json  # Fingerprint metadata
\end{verbatim}

\textbf{Critical Design Decision:} Layer 1 (features) uses filenames 
\textit{without} fingerprints because features are configuration-independent, they 
remain valid regardless of model hyperparameters or feature selection settings. 
Layer 2 (models) includes fingerprints in filenames because models become invalid 
when training configurations change. This naming convention prevents accidental 
cache misuse and enables visual inspection of cache validity.

\subsection{Cache Hit/Miss Logic}

The caching workflow follows this decision process:

\begin{enumerate}
\item Generate fingerprint from current configuration
\item Check if cached file exists with matching fingerprint
\item If \textbf{HIT}: Load cached data, record time saved
\item If \textbf{MISS}: Compute data, save to cache, record computation time
\end{enumerate}

latex\subsection{Current Implementation: Disk-Based Caching}


\begin{center}
\colorbox{yellow!20}{%
\begin{minipage}{0.9\textwidth}
\textbf{Working Draft Note:} 
rewriting this section right now still accurate here but need to adjust chapter 4 accordingly
\end{minipage}
}
\end{center}

The implemented caching system uses persistent SSD storage for all trained models. 
When a model is requested, the system checks for the existence of a cache file 
matching the configuration fingerprint and held-out subject identifier. If found, 
the model is loaded from disk; otherwise, it is trained and saved for future use.

\textbf{Lookup Strategy:}
\begin{enumerate}
\item Generate fingerprint from current configuration
\item Check if \verb|{fingerprint}_{subject}.joblib| exists on disk
\item If \textbf{HIT}: Load model from disk ($\sim$50ms)
\item If \textbf{MISS}: Train model ($\sim$300ms), save to disk, return
\end{enumerate}

This single-tier approach is well-suited to the experimental workflow of this 
thesis, where each model is typically accessed once per experimental run. The 
measured 4.5$\times$ speedup (Section~\ref{sec:verification}) demonstrates that 
disk-based caching alone provides substantial performance benefits for 
batch-oriented ML experimentation.





\begin{center}
\colorbox{yellow!20}{%
\begin{minipage}{0.9\textwidth}
\textbf{Working Draft Note:} 
delete here and put into future work?
\end{minipage}
}
\end{center}

\subsection{Future Extension: Two-Tier RAM+Disk Caching}

For production environments requiring interactive latency (e.g., real-time sleep 
analysis dashboards where the same subjects are analyzed repeatedly within minutes), 
the architecture could be extended with an in-memory cache layer:

\textbf{Two-Tier Lookup Strategy:}
\begin{enumerate}
\item \textbf{RAM Cache:} Check least-recently-used (LRU) cache in memory 
      ($\sim$0.1ms load time)
\item \textbf{Disk Cache:} If RAM miss, check persistent SSD storage ($\sim$50ms), 
      then store in RAM
\item \textbf{Compute:} If both miss, train model ($\sim$300ms), save to disk 
      and RAM
\end{enumerate}

\textbf{When Two-Tier Caching Provides Value:}
\begin{itemize}
\item \textbf{Interactive applications:} Dashboards where users repeatedly analyze 
      the same subjects with parameter adjustments
\item \textbf{High query frequency:} Systems processing $>$10 requests per minute 
      for the same models
\item \textbf{Low-latency requirements:} Applications requiring $<$100ms response times
\end{itemize}

\textbf{Why Not Implemented in This Thesis:}
The current experimental workflow accesses each cached model exactly once per run 
(128 LOSO folds $\times$ 1 access = 128 sequential lookups). Since there is no 
repeated access to the same model within a single experimental run, the RAM cache 
would never achieve hits beyond the initial disk load. Additionally, the 128 
cached models occupy $\sim$175 MB on disk but would require $\sim$700 MB of RAM 
if all were held in memory simultaneously, which could impact system stability 
during concurrent processing tasks.

For the batch-oriented experimentation conducted in this research, disk-based 
caching provides optimal cost-benefit trade-off: substantial speedup (4.5$\times$) 
with minimal memory overhead and complete persistence across sessions.

\textbf{Implementation Considerations for Future Work:}
A production implementation could use Python's \verb|functools.lru_cache| decorator 
or a dedicated caching library such as \verb|cachetools| to manage the RAM tier. 
The cache size would be tuned based on available memory (typically 10--50 models 
for a system with 4--8 GB available RAM). Complete implementation details and 
performance projections are documented in the project repository.

\subsection{Handling Data Changes and Cache Invalidation}

A critical challenge in ML experiment caching is ensuring cache validity when underlying data changes. Our implementation addresses this through a multi-layered integrity system embedded in the cache metadata. Each cached entry stores a configuration fingerprint---a SHA-256 hash of all preprocessing parameters, alongside explicit version identifiers (\texttt{data\_version} and \texttt{feature\_version}) and source file metadata.

When raw data modifications occur, three invalidation strategies are available:

\begin{enumerate}
\item \textbf{Version Bumping:} Incrementing the \texttt{data\_version} parameter in the configuration file causes the fingerprint to change, automatically invalidating all cached entries without requiring manual intervention. This approach preserves the cache structure while ensuring fresh computation.

\item \textbf{Cache Purging:} For significant data changes or version migrations, the global cache directory can be deleted entirely, triggering a complete recomputation on the next execution.

\item \textbf{Strict Validation:} An optional \texttt{strict\_validation} mode computes file checksums during cache lookup, detecting any byte-level changes to source files. While more computationally expensive, this provides cryptographic guarantees of data integrity.
\end{enumerate}

Our implementation defaults to version-based validation (Option 1), balancing integrity assurance with computational efficiency. The fingerprint system ensures that configuration drift---such as modified filter parameters or channel selections---automatically triggers recomputation, while unchanged configurations benefit from cached results. This approach aligns with reproducibility requirements in scientific computing, where identical inputs must yield identical outputs.



\section{Experimental Setup}
\label{sec:experiment_setup}

\subsection{Technical Environment}

Table~\ref{tab:technical_env} specifies the software environment for reproducibility. All code is developed and tested with Python 3.12.2 and the following dependencies on an HP ProBook x360 435 G7 laptop.

\begin{table}[ht]
\centering
\begin{tabular}{|l|l|l|}
\hline
\textbf{Component} & \textbf{Installed Version} & \textbf{Purpose} \\
\hline
Python & 3.12.2 & Core language \\
\hline
scikit-learn & 1.4.x+ & ML models (Random Forest, feature selection) \\
\hline
XGBoost & 3.0.2 & Gradient boosting classifier (primary model) \\
\hline
PyTorch & 2.7.0 & Deep learning framework for FNN model \\
\hline
MNE-Python & 1.9.0 & EEG data loading and signal processing \\
\hline
SciPy & 1.15.3 & Signal filtering, spectral analysis \\
\hline
NumPy & 2.1.3 & Numerical computations, array operations \\
\hline
Pandas & 2.2.3 & Data manipulation and organization \\
\hline
joblib & 1.5.1 & Parallel processing, model serialization \\
\hline
pickle & Built-in & Model/data serialization for caching \\
\hline
\textbf{Random Seed} & \textbf{42} & \textbf{All stochastic operations} \\
\hline
\end{tabular}
\caption{Technical environment and software versions for Sleep-EDF ML pipeline.}
\label{tab:technical_env}
\end{table}

\textbf{Hardware Environment:}

\begin{table}[ht]
\centering
\begin{tabular}{|l|l|}
\hline
\textbf{Component} & \textbf{Specification} \\
\hline
Machine & HP ProBook x360 435 G7 (Convertible Laptop) \\
\hline
Processor & AMD Ryzen 7 4700U with Radeon Graphics \\
 & (8 cores, 8 logical processors, 2000 MHz) \\
\hline
RAM (Installed) & 16.0 GB \\
\hline
RAM (Available during operation) & 1.6 GB \\
\hline
GPU & AMD Radeon Graphics (integrated) \\
 & No dedicated NVIDIA GPU \\
\hline
Storage & SSD on C: drive (primary boot device) \\
\hline
Operating System & Microsoft Windows 11 Pro (Build 26100) \\
\hline
System Type & x64-based PC with UEFI BIOS \\
\hline
\end{tabular}
\caption{Hardware environment specifications for Sleep-EDF ML pipeline development.}
\label{tab:hardware_env}
\end{table}

\textbf{Memory Management Strategy:}
With only 1.6 GB available RAM during operation, this pipeline employs aggressive memory optimization:
\begin{itemize}
	\item \textbf{Batch Processing:} Processes subjects in batches of 5--10 to keep feature matrices in memory
	\item \textbf{Memory Mapping:} Uses \texttt{numpy.memmap} for large arrays (loads only required portions into RAM on-demand)
	\item \textbf{Feature Caching:} Caches all 149 features to disk to avoid recomputation, reducing memory footprint of intermediate representations
	\item \textbf{Garbage Collection:} Explicit cleanup between subjects to free memory before next batch
	\item \textbf{Virtual Memory:} Leverages 15.4 GB page file if necessary, with acceptable I/O overhead
\end{itemize}

\textbf{Installation and Reproducibility:}
All dependencies are specified in the provided \texttt{requirements.txt} file. To replicate the development environment, run:

\begin{lstlisting}[language=bash]
pip install -r requirements.txt
\end{lstlisting}

\textbf{Performance Expectations:} This pipeline was developed and tested on a mobile workstation with 8 CPU cores and highly constrained available RAM (1.6 GB). Full LOSO cross-validation (128 folds) with feature caching requires approximately 12--16 hours of computation, with significant disk I/O overhead due to memory constraints. Results and timing metrics reported in Chapter~\ref{chap:results} are specific to this resource-constrained configuration. Users with more generous RAM availability will observe substantially faster execution, particularly for batch processing and model training phases.

\textbf{Note on Versions:} All installed versions are compatible and have been validated for this pipeline. For strict reproducibility in publication or deployment, consider pinning to exact versions in \texttt{requirements.txt}.
\subsection{Experimental Grid}

The experiments evaluate 18 configurations, as shown in Table~\ref{tab:exp_grid}.

\begin{table}[ht]
\centering
\begin{tabular}{|l|l|l|}
\hline
\textbf{Dimension} & \textbf{Options} & \textbf{Count} \\
\hline
Models & XGBoost, Random Forest & 2 \\
\hline
Correlation Thresholds & 0.75, 0.90, None & 3 \\
\hline
Top-K Features & 30, 50, None (all 149) & 3 \\
\hline
\textbf{Total Configurations} & 2 $\times$ 3 $\times$ 3 & \textbf{18} \\
\hline
\end{tabular}
\caption{Experimental grid: 18 configurations to evaluate.}
\label{tab:exp_grid}
\end{table}

Each configuration undergoes Leave-One-Subject-Out (LOSO) cross-validation with 128 folds, resulting in 2,304 total model training iterations.

\textbf{Note on FNN:} A Feedforward Neural Network (input$\rightarrow$64$\rightarrow$32$\rightarrow$5) 
is implemented and verified but excluded from final thesis experiments due to 
PyTorch dependency not being installed in the test environment. Preliminary 
results on 5 subjects achieved $\sim$78\% accuracy with successful caching 
(4.2$\times$ speedup on re-runs). The implementation is fully functional and 
could be included in future work with minimal additional effort ($<$1 hour 
computation time for full 18-configuration experiment once PyTorch is installed).




\begin{center}
\colorbox{yellow!20}{%
\begin{minipage}{0.9\textwidth}
\textbf{Working Draft Note:} 
as fnn is already partly implemented, but takes a lot of time to train (no dedicated gpu)fnn will not be used as of right now. maybe I am upgrading my pc soon to a laptop with dedicated gpu then i would like to incorporate the fnn model too
\end{minipage}
}
\end{center}


\subsection{Implementation Challenges and Resolutions}
\label{sec:implementation_challenges}


\begin{center}
\colorbox{yellow!20}{%
\begin{minipage}{0.9\textwidth}
\textbf{Working Draft Note:} 
note to myself, ask supervisor if these chapters should stay /or be canceld out - currently not sure about this chapters \end{minipage}
}
\end{center}

During development and testing, two critical bugs were identified and resolved. 
Documenting these issues provides insight into the subtleties of implementing 
ML experiment caching and may benefit researchers developing similar systems.

\subsubsection{Feature Mismatch Bug: Incomplete Fingerprinting}

\textbf{Symptom:} Models trained with 37 features crashed during loading when 
feature selection produced 33 features, with error: \texttt{X has 33 features, 
but XGBClassifier is expecting 37 features}.

\textbf{Root Cause:} The fingerprint generation included feature selection 
\textit{parameters} (correlation threshold, top-K count) but not the 
\textit{resulting feature names}. When different random seeds or data subsets 
produced different selected features under identical selection parameters, 
the fingerprint incorrectly indicated cache compatibility.

\subsubsection{Impact on Research Methodology}

These bugs were discovered through systematic testing before large-scale 
experiments, preventing potential corruption of results. The debugging process 
highlighted the importance of:
\begin{enumerate}
\item \textbf{Early verification testing} before committing to full experimental runs
\item \textbf{Explicit fingerprint design} that captures all sources of model variation
\item \textbf{Serialization-aware model design} for frameworks intended to support caching
\end{enumerate}

Both fixes are validated through the comprehensive test suite described in 
Section~\ref{sec:verification}.

\subsection{Cross-Validation Strategy}

\textbf{Leave-One-Subject-Out (LOSO)} cross-validation is the primary evaluation strategy:
\begin{itemize}
\item 128 folds (one per subject)
\item Each fold: 127 subjects for training, 1 subject for testing
\item Ensures true cross-subject generalization (no data leakage)
\item Total model trainings: 18 configurations $\times$ 128 folds = 2,304 trainings
\end{itemize}

\subsection{Caching Evaluation Scenarios}

Five scenarios will be executed to evaluate caching performance:

\begin{enumerate}
\item \textbf{First Run (Cold Start):} All cache misses, establishes baseline timing
\item \textbf{Identical Re-run:} All cache hits, measures maximum speedup
\item \textbf{Changed Preprocessing:} Invalidates Stage 1 and Stage 2 caches
\item \textbf{Changed Feature Selection:} Stage 1 and Stage 2 hit (selection at runtime)
\item \textbf{Add New Subjects:} Existing subjects hit, new subjects miss
\end{enumerate}

\section{Cache Verification Testing}
\label{sec:verification}

Prior to conducting the full 18-configuration experiment, the caching system 
underwent comprehensive verification to ensure correctness and measure performance.

\subsection{Test Design}


\begin{center}
\colorbox{yellow!20}{%
\begin{minipage}{0.9\textwidth}
\textbf{Working Draft Note:} 
maybe in final version cancel? as we use our real full grid + metrics and old test are not important anymore?
\end{minipage}
}
\end{center}

The verification protocol tested cache behavior across multiple scenarios using 
a reduced experimental grid (3 subjects $\times$ 8 configurations = 24 cache 
operations):

\begin{itemize}
\item \textbf{Subjects:} sub-1, sub-2, sub-3 (representative sample)
\item \textbf{Models:} XGBoost, Random Forest
\item \textbf{Feature Counts:} 30, 50, 149 (all)
\item \textbf{Correlation Thresholds:} 0.75, None
\item \textbf{Total Combinations:} $2 \times 3 \times 2 = 12$ configurations 
      (but $12 \times 2 = 24$ when counting both feature selection variants)
\end{itemize}

\textbf{Test Structure:}
\begin{enumerate}
\item \textbf{Cold Run:} Execute all 24 configurations with empty cache (expects 24 misses)
\item \textbf{Warm Run:} Re-execute identical configurations (expects 24 hits)
\item \textbf{Verification:} Compare predictions and metrics between runs
\end{enumerate}

\subsection{Verification Results}

Table~\ref{tab:cache_verification} shows the measured cache performance.

\begin{table}[ht]
\centering
\begin{tabular}{|l|c|c|c|c|c|}
\hline
\textbf{Run} & \textbf{Time} & \textbf{Hits} & \textbf{Misses} & \textbf{Hit Rate} & \textbf{Speedup} \\
\hline
Cold Start & 7.7s & 0 & 24 & 0\% & 1.0$\times$ (baseline) \\
\hline
Warm Re-run & 1.7s & 24 & 0 & 100\% & 4.5$\times$ \\
\hline
\end{tabular}
\caption{Cache verification results across 24 model training operations. Warm 
run achieved perfect cache hit rate with 4.5$\times$ speedup.}
\label{tab:cache_verification}
\end{table}

\subsection{Correctness Validation}

Beyond performance, the test suite verified cache correctness:

\begin{itemize}
\item \textbf{Prediction Identity:} Cached models produced bitwise-identical  predictions to freshly trained models (verified via \verb|numpy.array_equal()|)
\item \textbf{Metric Reproducibility:} Accuracy, F1, and kappa scores matched 
      to floating-point precision ($<10^{-10}$ difference)
\item \textbf{Zero Corruption:} All 24 pickle save/load cycles completed without errors
\item \textbf{Fingerprint Uniqueness:} Different configurations generated distinct 
      fingerprints (no hash collisions)
\end{itemize}

\subsection{Invalidation Testing}

Cache invalidation was verified by modifying configurations:

\begin{itemize}
\item \textbf{Changed Hyperparameter:} Modifying \verb|max_depth| from 6 to 8 
      $\rightarrow$ fingerprint changed $\rightarrow$ cache miss (correct)
\item \textbf{Changed Feature Selection:} Changing top-K from 30 to 50 $\rightarrow$ 
      fingerprint changed $\rightarrow$ cache miss (correct)
\item \textbf{Changed Held-Out Subject:} Testing sub-1 vs.\ sub-2 with identical 
      parameters $\rightarrow$ different fingerprints $\rightarrow$ no cross-fold 
      contamination (correct)
\end{itemize}

\subsection{Implications for Full Experiment}

These verification results establish confidence that:
\begin{enumerate}
\item The caching system correctly identifies cache hits/misses
\item Cached models are functionally identical to freshly trained models
\item Fingerprinting prevents inappropriate cache reuse across configurations
\item The measured 4.5$\times$ speedup on 24 operations extrapolates to substantial 
      time savings across the full 2,304-operation experiment (18 configs $\times$ 
      128 LOSO folds)
\end{enumerate}

The full experimental results (Chapter 5) will report cache performance across 
the complete dataset and experimental grid.


\section{Evaluation Plan}



\begin{center}
\colorbox{yellow!20}{%
\begin{minipage}{0.9\textwidth}
\textbf{Working Draft Note:} 
eventuell metriken leicht anpassen
\end{minipage}
}
\end{center}
\label{sec:evaluation_plan}

\subsection{Primary Metrics: Caching Efficiency}

The primary thesis contribution is evaluated through caching metrics:

\begin{itemize}
\item \textbf{Cache Hit Rate:} $\frac{\text{hits}}{\text{hits} + \text{misses}}$ per stage
\item \textbf{Time Saved:} $\sum(\verb|expected_compute_time| - \text{cache\_load\_time})$
\item \textbf{Speedup Factor:} $\frac{\text{cold\_start\_time}}{\text{cached\_run\_time}}$
\item \textbf{Storage Overhead:} Total cache size in MB/GB
\end{itemize}

\subsection{Secondary Metrics: Classification Performance}

Classification metrics validate that the caching framework produces correct results:

\begin{itemize}
\item \textbf{Accuracy:} Overall correct classifications
\item \textbf{Macro F1-Score:} Mean F1 across all classes (handles imbalance)
\item \textbf{Cohen's Kappa:} Agreement beyond chance
\item \textbf{Per-Class F1:} Identifies problematic classes (especially N1)
\item \textbf{Confusion Matrix:} 5$\times$5 aggregated across all folds
\end{itemize}

\subsection{Reproducibility Validation}

Reproducibility will be verified through:

\begin{itemize}
\item \textbf{Fingerprint Determinism:} Same config $\rightarrow$ same fingerprint (verified across multiple runs)
\item \textbf{Result Consistency:} Same fingerprint $\rightarrow$ identical metrics (within floating-point tolerance)
\item \textbf{Resumability Test:} Interrupt pipeline $\rightarrow$ resume from cache $\rightarrow$ identical final results
\end{itemize}

% ============================================================================
% CHAPTER 5: EXPECTED RESULTS (NEW CHAPTER FOR PROPOSAL)
% ============================================================================
\chapter{Expected Results}

This chapter outlines the anticipated findings and success criteria for evaluating the proposed intelligent caching framework. As this is a proposal, the values presented are educated estimates based on the system architecture and preliminary analysis. Actual measurements will be collected during the implementation phase and reported in the final thesis.

\section{Expected Caching Performance}

\subsection{Cache Hit Rates}

Based on the feature caching design, we expect the cache hit rates shown in Table~\ref{tab:expected_hit_rates}.

\begin{table}[ht]
\centering
\begin{tabular}{|l|c|l|}
\hline
\textbf{Scenario} & \textbf{Feature Cache Hit Rate} & \textbf{Explanation} \\
\hline
First Run (Cold Start) & 0\% & Features computed and cached \\
\hline
Identical Re-run & 100\% & Fingerprint matches \\
\hline
Changed Preprocessing & 0\% & Fingerprint changes, recompute \\
\hline
Changed Feature Selection Only & 100\% & Selection applied at query time \\
\hline
Add 20 New Subjects & 86\% & 128/148 subjects already cached \textcolor{red}{HIGHLIGHT FEATURES PERSISTENT !!! maybe delete here}  \\
\hline
\textbf{Typical Research Workflow} & \textbf{$>$90\%} & Model/hyperparam changes reuse cache \\
\hline
\end{tabular}
\caption{Expected feature cache hit rates across different experimental scenarios.}
\label{tab:expected_hit_rates}
\end{table}
\textcolor{red}{Also change 90 percent if 20 new subjects will be deleted }
The ``Typical Research Workflow'' row reflects that researchers often re-run experiments with minor parameter changes (e.g., different feature selection thresholds, different models), which should yield high cache hit rates for preprocessing and feature extraction.

\subsection{Measured Time Savings}

Table~\ref{tab:expected_time} presents the computation times and savings based on implementation benchmarks.

\begin{table}[ht]
\centering
\begin{tabular}{|l|c|c|c|}
\hline
\textbf{Operation} & \textbf{Cold Start} & \textbf{Cached} & \textbf{Speedup} \\
\hline
\multicolumn{4}{|c|}{\textit{Feature Extraction (Implemented -- Persistent Storage)}} \\
\hline
Preprocessing (per subject) & $\sim$5 sec & $\sim$5 sec & 1$\times$ (not cached) \\
\hline
Feature Extraction (per subject) & $\sim$25 sec & $\sim$0.1 sec & 250$\times$ \\
\hline
Full Pipeline (128 subjects) & $\sim$53 min & $\sim$14 sec & \textbf{227$\times$} \\
\hline
\multicolumn{4}{|c|}{\textit{LOSO Model Training (Proposed -- Intelligent Caching)}} \\
\hline
Single LOSO fold & $\sim$5 min & $\sim$1 sec & 300$\times$ \\
\hline
18-Config Experiment (2304 folds) & $\sim$16 hours & $\sim$4 min & \textbf{240$\times$} \\
\hline
\end{tabular}
\caption{Expected timing breakdown: Feature caching (implemented) vs. LOSO model caching (proposed). Feature extraction times are measured; LOSO model times are estimates based on preliminary benchmarks.}
\label{tab:expected_time}
\end{table}

\textbf{Key Insight:} The 227$\times$ speedup is achieved because the extraction of 149 EEG features per epoch (Stage 2) is cached once and reused across all 18 model configurations. This transforms a 16-hour experiment into a 4-minute experiment after the \textcolor{red}{complete} initial cold start. \textcolor{red}{check this section and think about keeoing feautre cache as time saved? better no}

\subsection{Measured Storage Requirements}

Table~\ref{tab:expected_storage} shows the actual storage overhead of the caching system using NPZ compression.

\begin{table}[ht]
\centering
\begin{tabular}{|l|c|c|}
\hline
\textbf{Cache Level} & \textbf{Per Subject} & \textbf{128 Subjects} \\
\hline
Stage 2 (149 Features, NPZ compressed) & $\sim$1.1 MB & $\sim$143 MB \\
\hline
\textbf{Total Cache} & $\sim$1.1 MB & \textbf{$\sim$143 MB} \\
\hline
\end{tabular}
\caption{Measured storage requirements for the feature caching system.}
\label{tab:expected_storage}
\end{table}



\begin{center}
\colorbox{yellow!20}{%
\begin{minipage}{0.9\textwidth}
\textbf{Working Draft Note:} 
Note outdated needs to be changed or deleted
\end{minipage}
}
\end{center}
\textcolor{red}{\textbf{Note:}} The current implementation caches only the extracted features (Stage 2). Preprocessed epochs (Stage 1) are computed on-the-fly during each run and not persisted to disk. This design prioritizes the feature extraction bottleneck, which accounts for $>$95\% of computation time. The compact NPZ format provides efficient compression using DEFLATE/zlib. For the 128-subject dataset with approximately 120,000 total epochs ($\sim$940 epochs per subject) and 195 features each, the uncompressed feature matrices would require approximately 183~MB (195 features $\times$ 120,000 epochs $\times$ 8 bytes = 187~MB). The compressed NPZ cache achieves $\sim$146~MB total, representing a compression ratio of approximately 1.25$\times$. This modest compression is typical for dense floating-point arrays, where the primary storage benefit comes from the efficient binary format rather than aggressive compression.

\section{Expected Classification Performance}


\begin{center}
\colorbox{yellow!20}{%
\begin{minipage}{0.9\textwidth}
\textbf{Working Draft Note:} 
should we keep this ? or weaken the benchmarks as this isnt our goal at all
\end{minipage}
}
\end{center}
While classification performance is secondary to the caching contribution, we expect results consistent with literature for EEG-based sleep staging.

\subsection{Overall Performance Expectations}

Based on the capstone project results and literature review, Table~\ref{tab:expected_classification} presents expected classification metrics.

\begin{table}[ht]
\centering
\begin{tabular}{|l|c|c|c|}
\hline
\textbf{Model} & \textbf{Accuracy} & \textbf{Macro F1} & \textbf{Cohen's Kappa} \\
\hline
XGBoost & 72--78\% & 58--68\% & 60--72\% \\
\hline
Random Forest & 68--75\% & 55--65\% & 55--68\% \\
\hline
\end{tabular}
\caption{Expected classification performance ranges based on LOSO cross-validation.}
\label{tab:expected_classification}
\end{table}

\subsection{Per-Class Performance Expectations}

Sleep stage classification exhibits known challenges with certain stages. Table~\ref{tab:expected_per_class} shows expected per-class F1 scores.

\begin{table}[ht]
\centering
\begin{tabular}{|l|c|l|}
\hline
\textbf{Stage} & \textbf{Expected F1} & \textbf{Notes} \\
\hline
Wake (W) & 75--85\% & Distinct EEG patterns \\
\hline
N1 (Light Sleep) & 20--40\% & Known difficulty---transitional stage \\
\hline
N2 (Sleep) & 75--85\% & Majority class, distinctive spindles \\
\hline
N3 (Deep Sleep) & 70--80\% & Strong delta waves \\
\hline
REM & 70--80\% & Distinctive mixed-frequency pattern \\
\hline
\end{tabular}
\caption{Expected per-class F1 scores with explanatory notes.}
\label{tab:expected_per_class}
\end{table}

\textbf{Note on N1 Classification:} The low expected F1 for N1 is consistent with literature and clinical practice. N1 is a brief transitional stage with EEG patterns similar to both Wake and N2, making it inherently difficult to classify. This is not a limitation of our approach but a known characteristic of the sleep staging problem.

\section{Expected Reproducibility Validation}

\subsection{Fingerprint Determinism}

We expect 100\% fingerprint determinism: the same configuration should always produce the same SHA-256 hash. This will be validated by running fingerprint generation multiple times with identical inputs.

\subsection{Result Consistency}

For identical fingerprints, we expect:
\begin{itemize}
\item Classification metrics identical within floating-point precision ($<$0.001\% variance)
\item Cached data byte-identical across loads (verified via checksum)
\item Model predictions identical for deterministic algorithms (XGBoost, RF with fixed seed)
\end{itemize}

\textbf{Note:} FNN models may exhibit minor variance ($\pm$0.5\%) due to GPU non-determinism, which will be documented as a known limitation.

\section{Success Criteria}




\begin{center}
\colorbox{yellow!20}{%
\begin{minipage}{0.9\textwidth}
\textbf{Working Draft Note:} 
Educated guesses so far:
\end{minipage}
}
\end{center}

The thesis will be considered successful if the following criteria are met:

\begin{enumerate}
\item \textbf{Caching Efficiency:}
  \begin{itemize}
  \item Cache hit rate = 100\% for repeated experiments with identical configurations
  \item Speedup factor $>$100$\times$ for cached features (achieved: 227$\times$)
  \item Storage overhead $<$200 MB for full 128-subject dataset (achieved: $\sim$143 MB)
    \item Speedup factor $>$10$\times$ for cached loso folds (128 subjects) (achieved: \textcolor{red}{placholder}$\times$)
  \end{itemize}

\item \textbf{Reproducibility:}
  \begin{itemize}
  \item 100\% fingerprint determinism
  \item Result variance $<$0.1\% for identical configurations
  \item Successful pipeline resumption after interruption
  \end{itemize}

\item \textbf{Classification Validity:}
  \begin{itemize}
  \item Overall accuracy $>$70\% (validates correct implementation)
  \item Macro F1 $>$55\% (handles class imbalance)
  \item Results consistent with published literature
  \end{itemize}
\end{enumerate}

% ============================================================================
% CHAPTER 6: SUMMARY AND FUTURE WORK
% ============================================================================
\chapter{Summary and Future Work}

\section{summary}
\textcolor{orange}{tba}

\section{Future Work}

Several extensions are planned or proposed for future research:

\begin{itemize}
\item \textbf{FNN Evaluation:} Complete the experimental evaluation of the implemented FNN model across all 18 configurations.
\item \textbf{Cross-Experiment Cache Reuse:} Enable cache sharing across different experiments with compatible preprocessing parameters.
\item \textbf{Distributed Caching:} Extend the framework to support multi-node caching for large-scale experiments.
\item \textbf{Integration with MLOps Tools:} Develop plugins for MLflow, DVC, or ZenML to bring fingerprint-based caching to existing workflows.
\item \textbf{Cross-Domain Validation:} Apply the framework to other ML domains (NLP, computer vision) to validate generalizability.
\item \textbf{Two level cache:} Last 10 saved on RAM, Everything on SSD see 4.1.6.
\end{itemize}
\subsection{Scalability Considerations for Large-Scale Experiments}

While this thesis addresses a controlled experimental scenario (128 subjects, 18 configurations, 2-stage caching), scaling the framework to production environments would introduce significant challenges that merit consideration.

\subsubsection{Cache Key Explosion}



\begin{center}
\colorbox{yellow!20}{%
\begin{minipage}{0.9\textwidth}
\textbf{Working Draft Note:} 
check the math and adjust if combination grid changes + autmatic invaldiation / deleting 
\end{minipage}
}
\end{center}

The current implementation manages a tractable configuration space. However, extending the framework to support dynamic dataset compositions would dramatically increase complexity. Consider the combinatorial explosion:

\begin{itemize}
\item \textbf{Current scope:} 1 fixed dataset $\times$ 18 configurations = 18 primary cache entries
\item \textbf{Dynamic subject combinations:} With 128 subjects, there exist $2^{128} \approx 3.4 \times 10^{38}$ possible subject subsets, clearly intractable for exhaustive caching
\item \textbf{Full configuration space:} Combining dynamic scopes ($\sim$500 realistic combinations) with preprocessing variants (48), CV strategies (3), models (4), feature strategies (6), hyperparameter grids ($\sim$10), and random seeds (3) yields approximately 259 million potential cache keys
\end{itemize}

This analysis suggests that naive ``cache everything'' strategies would not scale. Future implementations could explore selective caching policies that prioritize high-reuse, high-computation-cost artifacts while accepting recomputation for low-cost or rarely-accessed configurations.

\subsubsection{Dynamic Dataset Scope Management}

A key open question concerns how to reproducibly define and cache arbitrary subject combinations:

\begin{enumerate}
\item \textbf{Enumeration-based:} Predefined scopes (e.g., \texttt{dev-3}, \texttt{quality-high}, \texttt{full-128}) offer reproducibility but limited flexibility
\item \textbf{Hash-based identification:} Sorting subject IDs and hashing the resulting list provides deterministic keys for any combination, but does not address the storage explosion
\item \textbf{Progressive inclusion strategies:} Systematically adding subjects (1 $\rightarrow$ 2 $\rightarrow$ ... $\rightarrow$ 128) could enable incremental caching with $O(n)$ rather than $O(2^n)$ storage requirements
\end{enumerate}

\subsubsection{Hierarchical Invalidation at Scale}



\begin{center}
\colorbox{yellow!20}{%
\begin{minipage}{0.9\textwidth}
\textbf{Working Draft Note:} 
more a note section as finalisiert will be done after the piepline is complete and running and the implementation stands
\end{minipage}
}
\end{center}

The current two-stage invalidation model (preprocessing changes invalidate features; feature changes invalidate downstream results) could be extended to a more granular hierarchy:

\begin{enumerate}
\item \textbf{Per-subject epoching:} Stable, rarely invalidated
\item \textbf{Per-subject features:} Invalidated only on feature extraction code changes
\item \textbf{Aggregated datasets per scope:} Short TTL (e.g., 3 days) to prevent cache bloat
\item \textbf{Hyperparameter tuning intermediates:} Potentially excluded from caching entirely, as recomputation is typically fast
\item \textbf{Final models and results:} Persistent for reproducibility
\end{enumerate}

\subsubsection{Semantic Configuration Similarity}

An intriguing research direction concerns whether semantically similar configurations could share cached results. For example, if hyperparameter set $A$ produces results within 0.1\% of set $B$, could $B$'s cached results satisfy queries for $A$? This ``approximate caching'' concept parallels semantic caching in LLM applications but remains unexplored for ML experiment pipelines.

\subsubsection{Practical Implications}

These scalability considerations reinforce a key finding of this thesis: for domain-specific ML pipelines with controlled experimental designs, simple hierarchical caching with SHA-256 fingerprinting provides substantial benefits without requiring sophisticated ML-based cache prediction or semantic similarity matching. The complexity of such advanced techniques may only be justified in multi-user, multi-project environments where cache access patterns are less predictable and marginal efficiency gains translate to significant cost savings.



\chapter{Old versions - not yet deleted as "backup" within the document}

This thesis proposal presents an intelligent caching framework for optimizing ML experiments, addressing the critical challenges of computational efficiency and reproducibility in modern machine learning research.

\section{Research Questions Addressed}

The proposed framework addresses the central research question through four supporting aspects:

\begin{itemize}
\item \textbf{RQ1 (Efficiency):} The hierarchical caching system with Stage 2 (feature extraction) caching achieves 100\% cache hit rates and 227$\times$ speedup on repeated experiments.

\item \textbf{RQ2 (Reproducibility):} SHA-256 configuration fingerprinting ensures that identical configurations always produce identical cache keys, enabling reliable experiment resumption and result verification.

\item \textbf{RQ3 (Scalability):} The framework will be evaluated on both pilot (10 subjects) and full (128 subjects) datasets to demonstrate linear scaling of cache benefits.

\item \textbf{RQ4 (Trade-offs):} Storage overhead ($\sim$143 MB) is minimal compared to computational savings (16+ hours for full experimental grid).
\end{itemize}

\section{Expected Contributions}

This thesis is expected to make the following contributions:

\begin{enumerate}
\item \textbf{Theoretical Contribution:} A two-dimensional evaluation framework (WHY $\times$ WHERE) for analyzing caching strategies in ML pipelines.

\item \textbf{Practical Contribution:} A working implementation of fingerprint-based hierarchical caching that can be adapted to other ML workflows.

\item \textbf{Empirical Contribution:} Quantified efficiency gains and reproducibility validation through comprehensive experiments on the BOAS dataset.

\item \textbf{Methodological Contribution:} Guidelines for implementing LOSO-aware caching that prevents data leakage in cross-validation.
\end{enumerate}

\section{Limitations}

The following limitations are acknowledged:

\begin{itemize}
\item \textbf{Single Dataset:} Results are specific to BOAS; generalization to other domains requires further validation.
\item \textbf{Domain-Specific:} The 149-feature set is designed for EEG analysis; other applications would require different feature engineering.
\item \textbf{Hardware-Specific Timing:} Absolute timing results depend on the specific hardware configuration used.
\item \textbf{Feature Cache Only:} The current implementation caches extracted features (Stage 1--2). LOSO model caching (Stage 3) is designed but not implemented due to complexity; this is designated as future work.
\item \textbf{FNN Excluded:} The FNN model is implemented but excluded from thesis experiments due to time constraints.
\end{itemize}




% ============================================================================
% BACK MATTER
% ============================================================================
\backmatter%
\addtoToC{Bibliography}
\bibliographystyle{IEEEtranS}
\bibliography{references}
\end{document}